If assessing the quality of generated artifacts is important and no reference datasets or suitable comparison metrics exist, researchers \should use human validation for LLM outputs. If they do, they \must define the measured construct (e.g., usability, maintainability) and describe the measurement instrument in the \paper. Researchers \should consider human validation early in the study design, and not as an afterthought. They \should build on established reference models for human-LLM comparison. When developing or adapting measurement instruments, researchers \must share them. When LLMs replace, rather than augment, humans in research tasks such as annotating software artifacts, coding interview transcripts, or simulating study participants, researchers \must explain whether and how the replacement is justified and \should ground it with inter-model and model-to-human agreement. For qualitative coding of interpretive data, they \should justify why an LLM is methodologically appropriate. When aggregating LLM judgments, methods and rationale \should be reported and inter-rater agreement \should be assessed. Confounding factors \should be controlled for; where applicable, researchers may perform a power analysis to estimate the required sample size. For value-laden or culturally contingent constructs, researchers \should describe rater demographics beyond expertise and discuss potential demographic biases. When evaluating agentic tools, user feedback on the agent's proposed actions \should be assessed and acceptance statistics reported.